\subsection{分裂扩张}

定义2. --- 设K为一个域，且 \((f_{i})_{i\in I}\) 为 K{[}X{]} 中非常数多项式的一个族。所谓 \((f_{i})_{i\in I}\) 的一个分裂扩张，我们指的是K的任意满足以下性质的扩张E：

\begin{enumerate}
\def\labelenumi{\alph{enumi})}
\item
  对于每个 \(i \in \mathbf{I}\) ，多项式 \(f_i\) 在 \(\mathbf{E}[\mathbf{X}]\) 中分裂成若干个一次多项式的乘积。
\item
  对于每个 \(i \in I\) ，令 \(\mathbf{R}_i\) 为多项式 \(f_i\) 在 \(E\) 中的根集，则 \(E = K\left(\bigcup_{i \in I} \mathbf{R}_i\right)\) 。
\end{enumerate}

有时使用“分裂域”一词代替“分裂扩张”。

注记。--- 1) 对于每个 \(i \in I\) ，令 \(c_i\) 是 K 的一个非零元素，并令 \(f_i' = c_i f_i\) 。那么显然，族 \((f_i)_{i \in I}\) 的每个分裂扩张也是族 \((f_i')_{i \in I}\) 的一个分裂扩张，反之亦然。特别地，在研究分裂扩张时，我们可以仅限于考虑首一多项式的情形。

\begin{enumerate}
\def\labelenumi{\arabic{enumi})}
\setcounter{enumi}{1}
\tightlist
\item
  假设I是有限的，并设 \(f = \prod_{i \in I} f_i\) 。利用多项式分解为不可约因子的唯一性
\end{enumerate}

（于 E{[}X{]} 中）（IV，第 13 页，命题 13），我们可以容易地证明，多项式 f 的一个分裂扩张是族 \((f_{i})_{i \in I}\) 的一个分裂扩张，反之亦然。换言之，有限族的情形可以归结为单个多项式的情形。

\begin{enumerate}
\def\labelenumi{\arabic{enumi})}
\setcounter{enumi}{2}
\tightlist
\item
  设 \(f \in K[X]\) 是一个次数为 \(\geqslant 1\) 的多项式，并设 E 是 f 的一个分裂扩张。若 \(x_{1}, \ldots, x_{n}\) 是 f 在 E 中的根，则我们有 \(\mathrm{E} = \mathrm{K}(x_{1}, \ldots, x_{n})\) ，且 \([E : K]\) 是有限的（V，第18页，定理2）；但E可能不同于由单个根生成的子域 \(\mathrm{K}(x_{1}), \ldots, \mathrm{K}(x_{n})\) ；即使f是不可约的，这种情况也可能发生 \(^{1}\) 。然而我们注意到，当f不可约时，这些域 \(\mathrm{K}(x_{i})\) 在 K 上的次数均为 n，并且当 E 等于其中之一时，我们有 \([E : K] = n\) ，因此 \(\mathrm{E} = \mathrm{K}(x_{1}) = \cdots = \mathrm{K}(x_{n})\) 。
\end{enumerate}

命题 4. --- 设 K 为一个域，且 \((f_{i})_{i\in I}\) 为 K{[}X{]} 中非常数多项式的一个族；那么存在族 \((f_{i})_{i\in I}\) 的一个分裂扩张。

我们可以取这些多项式 \(f_{i}\) 为首一的（注记1）。设 \(i \in I\) 并令 \(f_{i}\) 的次数为 \(d_{i}\) 。根据IV，第73页，命题5，存在一个交换代数 \(A_{i}\) （在 \(K\) 上），它不约化为 0，以及元素 \(\xi_{i,1}, \ldots, \xi_{i,d_{i}}\) （属于 \(A_{i}\)），使得：

\begin{enumerate}
\def\labelenumi{\alph{enumi})}
\item
  该代数 \(\mathbf{A}_i\) 由 \((\xi_{i,1},\dots,\xi_{i,d_i})\) 生成；
\item
  我们有 \(f_{i}(\mathbf{X}) = \prod_{k=1}^{d_{i}} (\mathbf{X} - \xi_{i,k})\) 在 \(\mathbf{A}_i[\mathbf{X}]\) 中成立。
\end{enumerate}

设 A 为代数族 \((\mathbf{A}_{i})_{i\in I}\) 的张量积，并且令 \(\varphi_{i}\) 为 \(A_{i}\) 到 A 中的典范同态（III，第 470 页）。此时代数 A 是交换的且不约化为 0；因此由 Krull 定理（I，第 104 页），在 A 中存在一个极大理想 a，并且 \(E = A/a\) 是域K的扩张。

用 \(\psi\) 表示 A 到 E 的典范同态，并令 \(x_{i,k} = \psi (\varphi_i(\xi_{i,k}))\) （对 \(i\in I\) 及 \(1\leqslant k\leqslant d_i\)）。由于代数 A 由 \(\cup \varphi_i(A_i)\) 生成，扩张 E 由族 \((x_{i,k})\) 生成。进一步地，对每个 \(i\in I\) 我们有

\(f_{i}(\mathbf{X}) = \prod_{k=1}^{d_{i}} (\mathbf{X} - x_{i,k})\) 在 E{[}X{]} 中成立。因此，E 是族 \((f_{i})_{i \in I}\) 的一个分裂扩张。

命题5. --- 设 K 是一个域， \((f_i)_{i \in I}\) 为 K{[}X{]} 中非常数多项式的一个族，E 是 K 的一个扩张，且 F、F' 是 E 的子扩张，它们各自是 \((f_i)_{i \in I}\) 的一个分裂扩张。则 F = F'。

设 \(\mathbf{R}_i\) 为多项式 \(f_i\) 在 \(\mathbf{E}\) 中的根集，且 \(\mathbf{R} = \bigcup_{i\in I}\mathbf{R}_i\) 。由于 \(f_i\) 是

F{[}X{]} 中若干一次多项式的乘积，我们有 \(R_{i} \subset F\) 。根据定义2，我们有 \(F = K(R)\) ；同样地，我们发现 \(F' = K(R)\) 。

推论。--- 设 K 是一个域， \((f_{i})_{i\in I}\) 为 K{[}X{]} 中非常数多项式的一个族，且 F, F' 是 \((f_{i})_{i\in I}\) 的分裂扩张。则存在一个从 F 到 F' 的 K-同构。

这由命题5以及V，第13页，命题4的推论得出。

